% Two equivalent coordinate views of the four categorical FDE values.
\begin{figure}[t]
\centering
\begin{tikzpicture}[
  x=1.65cm,
  y=1.65cm,
  >=Latex,
  line cap=round,
  line join=round,
  every node/.style={font=\small}
]
  % Raw support square.
  \begin{scope}[shift={(0,0)}]
    \draw[->,thick] (-0.08,0) -- (1.28,0);
    \draw[->,thick] (0,-0.08) -- (0,1.28);
    \draw[black!45] (0,0) rectangle (1,1);

    \filldraw[fill=white,thick] (0,0) circle (2.2pt)
      node[below left] {$\Nval=0$};
    \filldraw[fill=white,thick] (1,0) circle (2.2pt)
      node[below right] {$\Tval=1$};
    \filldraw[fill=white,thick] (0,1) circle (2.2pt)
      node[above left] {$\Fval=i$};
    \filldraw[fill=white,thick] (1,1) circle (2.2pt)
      node[above right] {$\Bval=1+i$};

    \node[font=\bfseries] at (0.52,1.48) {support coordinates};
    \node at (0.52,-0.32) {$z_{\mathrm{s}}=v^+ + i v^-$};
  \end{scope}

  % Coordinate-change arrow.
  \draw[-{Latex[length=2.8mm]},very thick,blue!60!black]
    (1.45,0.55) -- (2.55,0.55);
  \node[align=center,text=blue!55!black,font=\scriptsize] at (2.00,0.28)
    {affine change};

  % Rotated truth/information diamond.
  \begin{scope}[shift={(3.85,0.55)}]
    \draw[->,thick] (-1.30,0) -- (1.34,0);
    \draw[->,thick] (0,-1.28) -- (0,1.34);
    \draw[black!45] (0,-1) -- (1,0) -- (0,1) -- (-1,0) -- cycle;

    \filldraw[fill=white,thick] (1,0) circle (2.2pt)
      node[above right] {$\Tval=1$};
    \filldraw[fill=white,thick] (-1,0) circle (2.2pt)
      node[above left] {$\Fval=-1$};
    \filldraw[fill=white,thick] (0,1) circle (2.2pt)
      node[above right] {$\Bval=i$};
    \filldraw[fill=white,thick] (0,-1) circle (2.2pt)
      node[below right] {$\Nval=-i$};

    \node[font=\bfseries] at (0,1.48) {truth/information coordinates};
    \node at (0,-1.38) {$z_{\mathrm{ti}}=(v^+-v^-)+i(v^++v^--1)$};
  \end{scope}
\end{tikzpicture}
\caption{Two equivalent embeddings of the four categorical FDE values. The
raw coordinate preserves the positive/negative support square. The transformed
coordinate separates classical truth polarity on the real axis from
glut/gap polarity on the imaginary axis. The figure represents a verified
coordinate transformation, not a new complex-valued logical connective.}
\label{fig:complex-coordinate-geometry}
\end{figure}
